% =============================================================================
% Introduction --- ZVF as per-step zero-gradient sentinel under group-relative
% advantages, with a within-task stratified audit of the pooled correlation
% Target: <= 1 page
% =============================================================================
\section{Introduction}
\label{sec:intro}

Group-relative policy optimization (GRPO) and its descendants now underlie a
large fraction of recent reasoning-RL pipelines~\citep{shao2024deepseekmath,%
guo2025deepseekr1nature}, replacing the value-function critic of
PPO~\citep{schulman2017proximal} with a leave-one-out advantage computed
within a small group of rollouts that share a prompt. Because the advantage
is the within-group reward minus the within-group mean (or median),
\emph{any group in which all rollouts receive identical rewards contributes a
zero advantage to every token} and therefore zero gradient. We call the
share of such groups in a training batch the \textbf{Zero-Variance Fraction}
(ZVF).

\paragraph{ZVF is a per-step sentinel, not a cross-task predictor.}
The framing we adopt in this paper is narrow and mechanical. At step~$t$,
$\mathrm{ZVF}_t = 1$ implies that every rollout group in the batch has zero
reward variance and therefore zero reward-gradient contribution; the
complementary $\mathrm{GU}_t = 1-\mathrm{ZVF}_t$ is an online \emph{utilization
meter} for the reward channel of the GRPO update. Pairing $\mathrm{ZVF}_t$
with batch reward mean separates all-wrong collapse ($\mathrm{ZVF}{\approx}1$,
mean reward $\approx 0$) from all-correct saturation ($\mathrm{ZVF}{\approx}1$,
mean reward $\approx 1$) from informative-contrast training
($\mathrm{ZVF}{<}1$, intermediate mean reward). This is the regime in which
ZVF is genuinely useful as instrumentation, and it is what caught the
sparse-reward tool-use failures in our benchmark, where every Tinker run
sat at $\mathrm{ZVF}{=}1$ from step~1 (Section~\ref{sec:zvf-analysis},
Figure~\ref{fig:zvf-sentinel}).

\paragraph{The pooled cross-task correlation is task identity in disguise.}
Earlier drafts of this paper, and several concurrent ``GRPO failure-mode''
analyses, summarized ZVF by pooling across tasks and computing a single
correlation between mean ZVF and final reward. Across the 15 cross-task
runs reported in earlier drafts, mean ZVF and last-10 reward correlate at
Pearson $r{=}-0.769$ ($p{=}0.0008$). On a larger $N{=}29$ sample with
released step-level traces, the pooled correlation drops to $r{=}-0.44$.
\emph{Both numbers are dominated by task identity, not by GRPO dynamics:}
tool-use runs sit at $\mathrm{ZVF}{\approx}1$ and reward $0$, GSM8K runs
sit at $\mathrm{ZVF}{\approx}0.085$ and reward $0.3$--$0.9$. Stratifying
within GSM8K, the correlation flips sign to $r{=}+0.40$ ($p{=}0.09$,
$N{=}19$); restricting to the seven GSM8K runs that vary model rather than
hyperparameter, $r{=}+0.13$ ($p{=}0.79$). The ``ZVF predicts performance''
story disappears under stratification and partial-correlation controls
(Section~\ref{sec:zvf-analysis}, Table~\ref{tab:evidence-grade}).

\paragraph{This paper is a focused statistical audit.}
We do not propose a new GRPO algorithm and we do not propose ZVF as a
benchmark scoring statistic. We make two paired claims and audit them. The
positive claim is that per-step ZVF inside a single task is a useful
zero-gradient sentinel and utilization meter; we provide a triage table
mapping (ZVF, reward) regimes to suggested actions. The negative claim is
that the pooled cross-task ZVF--reward correlation cited in earlier drafts
and concurrent work is confounded by task identity, does not survive
within-task stratification, and should not be used to rank GRPO runs.

This main-track paper is the focused sentinel/stratification claim in the larger
program. The adaptive controller and the GRPO--PPO--SAO unification are companion
proposals and do not upgrade the evidence reported here.

\paragraph{Contributions.}
(i)~A clean formalization of ZVF and its mechanical coupling to group size
and baseline accuracy under group-relative advantage estimators
(Section~\ref{sec:zvf-analysis} and Appendix~\ref{sec:appendix:zvf-formalization}).
(ii)~A positive sentinel demonstration on tool-use runs (Qwen3-32B and
Llama-3.1-8B, $\mathrm{ZVF}{=}1$, $\mathrm{GU}{=}0$, reward $0$ from step~1)
contrasted with a GSM8K run whose GU and reward both move, plus a triage
table mapping (regime, ZVF, reward) to action
(Section~\ref{sec:zvf-analysis}).
(iii)~A stratified audit of the pooled cross-task ZVF--reward correlation:
the pooled $r{=}-0.769$ collapses under within-task stratification on
GSM8K to a non-significant, sign-flipped $r{=}+0.40$ (and $r{=}+0.13$,
$N{=}7$ model-varying); we read this as a methodological warning, not a
defeat of the mechanical claim.
(iv)~A cross-stack benchmark of step-level ZVF traces on 79 GRPO runs across
7 RL libraries, 5 model families, and total parameter scales 0.6B--671B
(active 0.6B--37B for the largest mixture-of-experts checkpoints), released
as an artifact rather than headlined as a benchmark contribution
(Appendix~\ref{sec:run-manifest}).
(v)~Released step-level ZVF/GU traces, formalization, and reanalysis code.
